% Original vector diagram corresponding to source Fig.25.2, PDF111 / printed99.
% The two solid chi lines are distinct internal edges; both ell1,ell2 flow L->R.
% k=ell1+ell2. Dashed external phi stubs are displayed for orientation, but
% the reported iPi is amputated. S=2 exchanges the two internal chi lines.
% Requires basic TikZ only; no bitmap and no extra drawing library.
\begin{tikzpicture}[x=14mm,y=12mm,line width=.65pt,>=latex,
  font=\fontsize{10.5}{12}\selectfont]
  \draw[dashed] (-3,0) -- (-1.15,0);
  \draw[dashed] (1.15,0) -- (3,0);
  \draw (-1.15,0) .. controls (-.65,1.08) and (.65,1.08) .. (1.15,0);
  \draw (-1.15,0) .. controls (-.65,-1.08) and (.65,-1.08) .. (1.15,0);
  \fill (-1.15,0) circle[radius=.65mm];
  \fill (1.15,0) circle[radius=.65mm];
  \draw[->] (-2.6,.20) -- (-1.8,.20);
  \draw[->] (1.8,.20) -- (2.6,.20);
  \node[above] at (-2.2,.20) {$k$};
  \node[above] at (2.2,.20) {$k$};
  \node[below] at (-2.2,-.05) {$\varphi$};
  \node[below] at (2.2,-.05) {$\varphi$};
  \draw[->] (-.42,1.09) -- (.42,1.09);
  \node[above] at (0,1.09) {$\ell_1$};
  \draw[->] (-.42,-1.09) -- (.42,-1.09);
  \node[below] at (0,-1.09) {$\ell_2$};
  \node at (0,0) {$S=2$};
  \node at (0,-1.89) {截去外传播子后计算 $i\Pi(k^2)$};
\end{tikzpicture}
